\chapter{Etika \& Aturan Konseling AI}

\section{Prinsip Etika Konseling AI}
Luna AI beroperasi berdasarkan norma dan etika konseling digital yang mengutamakan keselamatan, kenyamanan, serta kerahasiaan pengguna.

\begin{lunacallout}{Prinsip Utama Layanan}
\begin{enumerate}
    \item \textbf{Non-Judgmental Space}: AI tidak memberikan penilaian moral terhadap emosi atau pengalaman pengguna.
    \item \textbf{Keramahan \& Empati}: Bahasa yang digunakan dirancang ramah, tenang, dan mendukung refleksi emosional.
    \item \textbf{Kerahasiaan Data}: Percakapan tidak dijual atau dibagikan kepada pihak ketiga non-medis tanpa izin.
\end{enumerate}
\end{lunacallout}

\section{Batasan Wewenang Luna AI}
Sangat penting bagi setiap pengguna untuk memahami bahwa Luna AI adalah sistem pendukung kecerdasan buatan, bukan psikolog atau psikiater manusia berlisensi.

\begin{lunawarning}[Batasan Medis (Medical Disclaimer)]
Luna AI \textbf{TIDAK BERWEWENANG} untuk:
\begin{itemize}
    \item Memberikan diagnosis medis resmi (misal: Klinis Depresi, Bipolar, Schizophrenia).
    \item Mengeluarkan resep obat-obatan psikiatri.
    \item Menggantikan sesi terapi langsung dengan profesional medis berlisensi.
\end{itemize}
\end{lunawarning}

\section{Tata Tertib Interaksi Pengguna}
Pengguna diharapkan menjaga etika interaksi saat mengoperasikan layanan Luna AI:
\begin{itemize}
    \item Dilarang menggunakan percakapan untuk mempromosikan tindakan ilegal atau membahayakan orang lain.
    \item Dilarang mencoba melakukan \textit{jailbreaking} atau eksploitasi prompt untuk memancing keluaran berbahaya dari LLM.
    \item Diharapkan memberikan umpan balik (\textit{feedback rating}) secara jujur untuk membantu peningkatan performa AI.
\end{itemize}

\appfigurepair[0.25\textwidth]{images/terms_page_1.jpeg}{Bagian 1: Etika \& Disclaimer}{images/terms_page_2.jpeg}{Bagian 2: Krisis \& Privasi}{Syarat, Ketentuan \& Disclaimer Etika}{Tampilan lengkap panduan etika konseling AI, batasan wewenang medis, dan protokol keselamatan.}
